% !TeX root = ../main.tex

\chapter{基于异构语义表征增强与去噪的交通知识问答方法}

\section{引言}

\section{交通知识问答异构语义建模}

\subsection{异构语义表征增强问题定义}

\subsection{文本语义与图知识的数据关联}

\subsection{异构信息噪声与语义冲突建模}

\section{知识图语义表征增强策略}

\subsection{实体—关系上下文编码}

\subsection{任务条件的子图语义增强}

\subsection{图知识选择与伪语义抑制}

\section{文本—知识表征对齐与去噪策略}

\subsection{文本与图知识的分布对齐}

\subsection{跨源语义一致性约束}

\subsection{冲突语义检测与信息去噪}

\section{异构表征融合与答案推断}

\subsection{文本与知识的增强表示}

\subsection{异构语义非线性融合}

\subsection{答案预测与联合优化}

\section{实验结果与分析}

\subsection{交通知识问答数据复用与异构扰动设置}

\subsection{定量分析}

\subsection{噪声与缺失条件下的鲁棒性分析}

\subsection{消融实验}

\subsection{定性分析}

\section{本章小结}
